\documentclass[11pt,a4paper]{article}

\usepackage[utf8]{inputenc}
\usepackage[T1]{fontenc}
\usepackage[spanish,es-noquoting,es-nolists,es-noshorthands]{babel}
\usepackage{tgpagella}
\usepackage{inconsolata}
\usepackage{microtype}
\usepackage[a4paper,top=2.2cm,bottom=2.2cm,left=2.4cm,right=2.4cm]{geometry}
\usepackage{xcolor}
\usepackage{titlesec}
\usepackage{enumitem}
\usepackage{booktabs}
\usepackage{tabularx}
\usepackage{array}
\usepackage{needspace}
\usepackage{lastpage}
\usepackage{fancyhdr}
\setlength{\headheight}{14pt}
\usepackage[most]{tcolorbox}
\usepackage[colorlinks=true,linkcolor=arcablue,urlcolor=arcablue,citecolor=arcablue]{hyperref}

% ---------------------------------------------------------------- estilo
\definecolor{arcablue}{HTML}{1F3A63}
\definecolor{arcagrey}{HTML}{5A6472}
\definecolor{arcalight}{HTML}{EEF2F7}
\definecolor{arcarule}{HTML}{C9D3E0}

\titleformat{\section}{\normalfont\large\bfseries\color{arcablue}}{\thesection.}{0.6em}{}
\titleformat{\subsection}{\normalfont\normalsize\bfseries}{\thesubsection}{0.6em}{}
\titlespacing{\section}{0pt}{1.3ex plus .4ex minus .2ex}{0.8ex}
\titlespacing{\subsection}{0pt}{1.2ex plus .3ex}{0.6ex}

\setlist[itemize]{leftmargin=1.2em,itemsep=0.15ex,topsep=0.4ex,parsep=0pt}
\setlist[enumerate]{leftmargin=1.5em,itemsep=0.25ex,topsep=0.4ex,parsep=0pt}

\pagestyle{fancy}
\fancyhf{}
\renewcommand{\headrulewidth}{0.4pt}
\renewcommand{\headrule}{\hbox to\headwidth{\color{arcarule}\leaders\hrule height \headrulewidth\hfill}}
\fancyhead[L]{\footnotesize\color{arcagrey}Modelos Generativos Profundos · MIA · ICAI}
\fancyhead[R]{\footnotesize\color{arcagrey}Práctica final: ARCA}
\fancyfoot[C]{\footnotesize\color{arcagrey}\thepage\ de \pageref{LastPage}}

\newcommand{\code}[1]{{\small\texttt{#1}}}
\newcommand{\file}[1]{{\small\texttt{\textcolor{arcablue}{#1}}}}

\newtcolorbox{minimos}[1]{
  enhanced, breakable, colback=arcalight, colframe=arcablue,
  boxrule=0pt, leftrule=2.5pt, arc=1pt, left=8pt, right=8pt, top=6pt, bottom=6pt,
  fonttitle=\bfseries\small, title={#1}, coltitle=arcablue,
  attach boxed title to top left={xshift=8pt,yshift=-2pt},
  boxed title style={colback=arcalight,boxrule=0pt,arc=0pt}
}

\newtcolorbox{nota}{
  enhanced, breakable, colback=white, colframe=arcarule,
  boxrule=0.6pt, arc=1pt, left=8pt, right=8pt, top=5pt, bottom=5pt
}

\setlength{\parindent}{0pt}
\setlength{\parskip}{0.55ex}

% ---------------------------------------------------------------- documento
\begin{document}

\begin{center}
  {\Large\bfseries\color{arcablue} ARCA: Agente con Razonamiento, Conocimiento y Acción}\\[0.5ex]
  {\large Práctica final}\\[1.2ex]
  {\small Modelos Generativos Profundos · Máster en Inteligencia Artificial · ICAI · Curso 2026-2027}\\[0.8ex]
  {\small Repositorio: \url{https://github.com/kendrickcetina/dgm-arca}}
\end{center}

\vspace{0.6ex}
{\color{arcarule}\hrule height 0.8pt}
\vspace{1.2ex}

\section{Qué vais a construir}

Durante el curso hemos visto cómo un modelo de lenguaje aprende a razonar, cómo se le enseña
a usar herramientas, cómo se le conecta a una base de conocimiento y cómo todo eso se junta
en un agente que actúa en bucle. En esta práctica lo construís vosotros, de principio a fin,
\textbf{sobre un tema que elegís vosotros}.

La llamamos ARCA porque cada letra es una de las piezas: \textbf{R}azonamiento,
\textbf{C}onocimiento, \textbf{A}cción, y el \textbf{A}gente que las une.

El tema es libre. Un asistente para preparar unas oposiciones, un agente que lee la normativa
de un ayuntamiento, un tutor de matemáticas, un copiloto fiscal para autónomos, un entrenador
de programación competitiva. Lo que queráis, con una condición: que resuelva un problema a
alguien de verdad, aunque ese alguien seáis vosotros. Si al terminar podéis enseñar el
repositorio en una entrevista de trabajo, habréis hecho bien el trabajo.

Sé que vais a programar con ayuda de agentes de código, y no os lo voy a prohibir.
Precisamente por eso lo que puntúa no es escribir código, sino \textbf{decidir qué construir,
saber si funciona y explicar por qué}. Un agente de código os escribe un bucle GRPO en dos
minutos; lo que no puede hacer es elegir la tarea verificable de vuestro dominio, diseñar el
verificador, construir el corpus ni interpretar por qué la recompensa se estanca.

\begin{nota}
\textbf{Tres condiciones que el tema debe cumplir.} (1) Alguna tarea con \textbf{respuesta
comprobable automáticamente}: un número, una fecha, un código que pasa tests, una etiqueta.
Sin verificador no hay recompensa y sin recompensa no hay refuerzo. (2) \textbf{Herramientas
con sentido} en ese dominio: algo que consultar, algo que calcular, algo que hacer.
(3) \textbf{Documentación real} que podáis recopilar. Sin corpus no hay RAG.
\end{nota}

\section{Equipos, ritmo y forma de trabajar}

\textbf{Equipos de tres personas.} Repositorio compartido, un solo entregable por equipo, y
en la defensa oral pregunto a cualquiera de los tres sobre cualquier parte del código.

La práctica se construye \textbf{por capas, y se prueba desde el primer día}. No esperéis a
tenerlo todo: cada fase termina en un endpoint que funciona por sí solo, así que podéis
levantar la API con la fase 1 hecha y las otras tres devolviendo «pendiente».

\begin{center}
\small
\begin{tabularx}{\textwidth}{@{}l l l X@{}}
\toprule
\textbf{Fase} & \textbf{Carpeta} & \textbf{Endpoint} & \textbf{Cuándo} \\
\midrule
0. Propuesta      & \file{docs/}     & ---                & Primera semana, bloqueante \\
1. Razonamiento   & \file{rlm/}      & \code{/reasoning}  & Empezamos el lunes 21 de septiembre \\
2. Herramientas   & \file{tool\_use/}& \code{/tools}      & A continuación \\
3. Conocimiento   & \file{rag/}      & \code{/rag}        & A continuación \\
4. Agente         & \file{agent/}    & \code{/agent}      & Cierra el proyecto \\
\bottomrule
\end{tabularx}
\end{center}

El orden no es casual: primero el modelo aprende a pensar, luego le damos manos, luego memoria
externa, y al final lo ponemos a trabajar solo usando todo lo anterior. En la fase 4
reutilizaréis el modelo de la fase 1, las herramientas de la fase 2 y el retriever de la
fase 3. Pensad en el conjunto desde el principio.

Las fechas exactas de cada entrega las fijamos en clase. Cada README de carpeta tiene el
detalle completo; este documento es el resumen.

\section{Fase 0: la propuesta}

Antes de escribir código, una página por equipo. No es burocracia: es la parte que más
determina vuestra nota, porque un tema bien elegido hace que todo lo demás fluya.

Plantilla en \file{docs/00\_propuesta.md}. Quiero leer: el tema y el usuario; diez preguntas
reales que ese usuario haría; \textbf{qué tarea es verificable y cómo la verificaríais};
qué herramientas harían falta; de dónde sale el corpus; y qué puede salir mal.

Si no tenéis tema, en \file{docs/temas\_ejemplo.md} hay trece desarrollados (farmacia, legal,
finanzas, desarrollo, social, marketing, ciencia) con su tarea verificable, sus herramientas
y su corpus. Son para inspirar, no para elegir de una lista.

\textbf{Hasta que la propuesta esté aprobada, no empecéis la fase 1.} Cuando la apruebe os
devolveré por escrito qué se evalúa exactamente en vuestro tema.

\section{Fase 1: Razonamiento (\file{rlm/})}

Convertir un modelo de lenguaje en un modelo que piensa: que genere su razonamiento entre
\code{<think>} y \code{</think>} y después la respuesta entre \code{<answer>} y
\code{</answer>}. Reproducís, a pequeña escala y sobre vuestro dominio, el camino de
DeepSeek-R1: arranque en frío con SFT y después refuerzo con recompensas verificables.

\begin{minimos}{Mínimos de la fase 1}
\begin{enumerate}
  \item \textbf{Dataset verificable propio} (cientos de problemas) y un \textbf{verificador
  determinista} en \file{verifier.py}, con tests de sus casos límite.
  \item \textbf{Destilación}: generar trazas con un modelo profesor, filtrarlas con el
  verificador y quedarse con las correctas. Reportad la tasa de aceptación.
  \item \textbf{SFT} con LoRA sobre las trazas verificadas.
  \item \textbf{GRPO} con al menos tres recompensas: formato, exactitud y una
  \textbf{propia de vuestro dominio} que tendréis que justificar.
  \item \textbf{Un paso de GRPO implementado a mano} (ventajas, ratio de políticas, recorte,
  KL) con sus tests en verde.
  \item \textbf{Evaluación}: pass@1 en test para modelo base, tras SFT y tras GRPO; curvas de
  recompensa y de longitud; análisis de cinco fallos concretos.
\end{enumerate}
\end{minimos}

\textbf{Ficheros.} Ya hechos: \file{rewards.py} (recompensas de formato y exactitud, con
tests), \file{data.py} (prompt de sistema de R1-Zero y carga de datos),
\file{generate\_problems.py} (generador de problemas de ejemplo, ejecutable),
\file{inference.py} (detrás del endpoint). Vuestros: \file{verifier.py}, \file{distill.py},
\file{train\_sft.py}, \file{train\_grpo.py}, \file{grpo\_step.py}, \file{evaluate.py}.

\textbf{¿De dónde salen cientos de problemas? No de anotarlos a mano.} Las trazas de
razonamiento las genera el modelo profesor y las filtra vuestro verificador. Y los pares
(enunciado, respuesta) se consiguen con un generador programático, minando datos
estructurados, o con benchmarks públicos del dominio. Está todo en
\file{docs/datasets.md}, con tamaños de referencia y los errores típicos.

\textbf{Antes de empezar}, corred el entrenamiento de prueba: \code{uv run arca-smoke}. Es
GRPO de verdad en miniatura sobre GSM8K, con las mismas recompensas que usaréis. En diez
minutos veréis la recompensa de formato subir. Detalles en \file{smoke/README.md} y
\file{docs/dgx.md}.

\section{Fase 2: Herramientas (\file{tool\_use/})}

Un modelo solo produce texto: no puede consultar una API ni ejecutar código. Tool use es
dárselo. El modelo genera una llamada en un formato reconocible, la aplicación la detecta, la
ejecuta y le devuelve el resultado. Los cinco pasos que vimos en clase, implementados a mano.

\begin{minimos}{Mínimos de la fase 2}
\begin{enumerate}
  \item \textbf{Tres herramientas reales} de vuestro dominio: una que consulte una API HTTP
  externa de verdad, una que calcule o ejecute, y una que \textbf{actúe} con efecto
  observable (escribir un fichero, crear un evento, enviar a un webhook de pruebas).
  \item \textbf{Definiciones como JSON Schema} generadas desde Pydantic, con los argumentos
  validados antes de ejecutar; un error de validación vuelve al modelo como observación.
  \item \textbf{El flujo de cinco pasos escrito por vosotros}, sin frameworks de agentes, con
  soporte para llamadas encadenadas (multi-step) y varias llamadas en una generación
  (paralelas).
  \item \textbf{Banco de pruebas de treinta casos} con la herramienta y los argumentos
  esperados, incluyendo casos que \textbf{no} necesitan herramienta. Métricas: precisión de
  selección y de argumentos. Comparad el modelo con y sin herramientas.
\end{enumerate}
\end{minimos}

\textbf{Ficheros.} Ya hechos: \file{registry.py} (convierte funciones tipadas en definiciones
y valida), \file{parser.py} (detecta \code{<tool\_call>}), \file{tools.py} (dos herramientas
de ejemplo), \file{evaluate.py} (métricas). Vuestros: vuestras herramientas en
\file{tools.py}, \file{executor.py} y el banco en \file{bench/}.

\textbf{En esta fase no se usa \code{smolagents}.} Quiero ver vuestro bucle. En la fase 4 sí
lo permito, como contraste.

\section{Fase 3: Conocimiento (\file{rag/})}

El conocimiento del modelo se congela en su fecha de corte, y meterle documentos enteros en el
contexto no funciona. RAG es la alternativa: trocear el corpus, indexarlo, y recuperar solo
los fragmentos relevantes en el momento de la pregunta.

\begin{minimos}{Mínimos de la fase 3}
\begin{enumerate}
  \item \textbf{Corpus real} de vuestro dominio con ingesta reproducible (del orden de
  decenas de documentos o varios cientos de fragmentos), y su licencia documentada.
  \item \textbf{Dos estrategias de troceado comparadas} con números: la recursiva y la
  semántica basada en embeddings de frase.
  \item \textbf{Embeddings instruction-aware} (Qwen3-Embedding), almacén vectorial,
  \textbf{BM25} y \textbf{recuperación híbrida} con un peso $\lambda$ ajustable.
  \item \textbf{Conjunto dorado de cincuenta preguntas} anotadas con el fragmento que las
  responde, y \textbf{Recall@k y MRR} para los tres retrievers. El $\lambda$ se elige con
  datos, no a ojo.
  \item \textbf{Respuestas con citas} a los fragmentos usados, y una comprobación automática
  de que las citas existen.
\end{enumerate}
\end{minimos}

\textbf{Ficheros.} Ya hechos: \file{ingest.py} (lee PDF, Markdown, texto, HTML),
\file{chunking.py} (troceado recursivo y estadísticas), \file{embed.py} (Qwen3-Embedding con
instrucción), \file{retriever.py} (BM25), \file{evaluate.py} (Recall@k y MRR),
\file{generate.py} (prompt aumentado y extracción de citas). Vuestros: el troceado semántico,
el retriever denso, el híbrido, el corpus en \file{corpus/} y el conjunto dorado en
\file{gold/}.

El trabajo de verdad de esta fase es el conjunto dorado. Se escribe leyendo el corpus, y es
lo que os permite decir qué retriever es mejor en vuestro dominio en lugar de suponerlo.

\needspace{16\baselineskip}
\section{Fase 4: Agente (\file{agent/})}

Todo junto en un bucle: piensa, actúa, observa, y vuelve a pensar hasta responder. Es el
patrón ReAct, y es donde se juntan las tres fases anteriores.

\begin{minimos}{Mínimos de la fase 4}
\begin{enumerate}
  \item \textbf{Bucle ReAct propio} con \code{final\_answer}, límite de pasos, timeouts por
  herramienta y recuperación de errores. El retriever de la fase 3 entra como la herramienta
  \code{search\_knowledge\_base}.
  \item \textbf{Cerebro intercambiable} y comparado: modelo base, vuestro modelo de la
  fase 1, y un modelo con razonamiento nativo. ¿Ayuda de verdad haber entrenado el modelo?
  \item \textbf{Memoria} multi-turno que se resume automáticamente al superar un presupuesto
  de tokens.
  \item \textbf{Banco de veinte tareas} que requieran encadenar al menos dos herramientas,
  con tasa de éxito, número medio de pasos y coste en tokens. Comparad vuestro agente JSON
  con un agente de código.
  \item \textbf{Traza visible}: cada ejecución devuelve todos los pasos, y una interfaz
  mínima los muestra.
\end{enumerate}
\end{minimos}

\textbf{Ficheros.} Ya hechos: \file{prompts.py} (prompt ReAct), \file{memory.py} (memoria con
presupuesto), \file{evaluate.py} (métricas por cerebro), \file{ui.py} (visor de trazas).
Vuestros: \file{react\_agent.py}, la política de memoria y el banco en \file{bench/}.

\section{Todo se prueba desde la API}

\textbf{La corrección se hace llamando a vuestra API.} No leo vuestros notebooks para saber
si funciona: lanzo peticiones. Por eso los contratos de entrada y salida están fijados en
\file{api/schemas.py}. Podéis añadir campos opcionales a las respuestas, pero no renombrar ni
quitar los que ya están.

\begin{center}
\small
\begin{tabularx}{\textwidth}{@{}l l X@{}}
\toprule
\textbf{Endpoint} & \textbf{Envío} & \textbf{Devuelve} \\
\midrule
\code{GET /health} & --- & Estado de cada fase: \code{ready}, \code{pending} o \code{error} \\
\code{POST /reasoning} & \code{question}, \code{expected\_answer?} & \code{thinking},
  \code{answer}, \code{has\_valid\_format}, veredicto del verificador \\
\code{POST /tools} & \code{query}, \code{max\_turns?} & \code{answer} y \code{tool\_calls[]}
  con nombre, argumentos, resultado y latencia \\
\code{POST /rag} & \code{question}, \code{top\_k?}, \code{retriever?} & \code{answer},
  \code{chunks[]} con id, fuente y puntuación, y \code{citations[]} \\
\code{POST /agent} & \code{task}, \code{max\_steps?}, \code{brain?} &
  \code{final\_answer} y \code{steps[]} con pensamiento, acción y observación \\
\bottomrule
\end{tabularx}
\end{center}

\textbf{Cómo funciona mientras desarrolláis.} La API arranca siempre, aunque no hayáis hecho
nada: una fase sin implementar responde \code{501} con su nombre y un mensaje, y
\code{/health} os dice qué está listo. Así entregáis por capas y probáis cada fase en cuanto
la terminéis.

\textbf{Un ejemplo de lo que espero ver.}

{\footnotesize
\begin{verbatim}
$ curl -s localhost:8000/health
{"status": "ok", "team": "equipo-07", "domain": "fiscalidad de autonomos",
 "phases": {"reasoning": {"status": "ready"}, "tools": {"status": "pending"}, ...}}

$ curl -s -X POST localhost:8000/reasoning -H 'Content-Type: application/json' \
       -d '{"question": "...", "expected_answer": "120"}'
{"thinking": "Sea C la capacidad del deposito...", "answer": "120",
 "has_valid_format": true, "tokens_generated": 214,
 "verifier": {"is_correct": true, "predicted": "120", "expected": "120"}}
\end{verbatim}
}
\vspace{-1.6ex}

\begin{nota}
\textbf{Para la corrección.} La API tiene que levantar con \code{docker compose up api} en
una máquina limpia. La exponéis con \href{https://ngrok.com/}{ngrok} (\code{ngrok http 8000})
y me pasáis la URL. Antes de pasármela, probad los cuatro endpoints desde \code{/docs} y
comprobad que \code{/health} los marca como \code{ready}. Rellenad también \code{ARCA\_TEAM}
y \code{ARCA\_DOMAIN} en vuestro \code{.env}: salen en \code{/health} y me dicen a quién
estoy corrigiendo. Detalles en \file{api/README.md}.
\end{nota}

\needspace{14\baselineskip}
\section{Cómo se evalúa}

Los cuatro criterios se aplican \textbf{a cada fase por separado}, y cada fase pesa lo mismo
en la nota final.

\begin{center}
\small
\begin{tabularx}{\textwidth}{@{}l c X@{}}
\toprule
\textbf{Criterio} & \textbf{Peso} & \textbf{Qué miro} \\
\midrule
Implementación correcta & 30\,\% & El endpoint responde y hace lo que dice. Lo compruebo con
  un script que lanza peticiones, algunas parecidas a vuestros bancos de pruebas y otras que
  no habéis visto. \\
Completitud & 30\,\% & Los mínimos de la fase están hechos. Cada pieza que falte resta. \\
Interpretación y explicación & 30\,\% & Curvas y métricas explicadas, no solo pintadas.
  Comparaciones y ablaciones. Análisis de fallos. Cuaderno de experimentos honesto. Y la
  \textbf{defensa oral}. \\
Limpieza y calidad & 10\,\% & Sin marcadores pendientes. \code{uv sync} reproduce el entorno,
  los tests pasan, Docker levanta, README de portfolio, sin secretos en el repositorio. \\
\bottomrule
\end{tabularx}
\end{center}

La diferencia entre un trabajo correcto y uno excelente está en el tercer criterio. Quiero
que sepáis por qué vuestra recompensa es como es, qué hace el recorte en GRPO, por qué
elegisteis ese $\lambda$ y qué pasó cuando le quitasteis el modelo de razonamiento al agente.
Si el código lo escribió un agente y no sabéis explicarlo, se nota.

Cada fase se entrega con su parte del cuaderno de experimentos
(\file{docs/EXPERIMENTS\_plantilla.md}), fechado: una entrada que diga «probamos esto, no
funcionó, creemos que por aquello» vale más que diez que digan «todo bien». La rúbrica
completa está en \file{docs/rubrica.md}.

\section{La entrega final}

Cuando terminéis las cuatro fases tendréis un sistema completo. Tratadlo como un producto:

\begin{itemize}
  \item Repositorio público con README de portfolio: qué hace, para quién, una demo, un
  diagrama, la tabla de métricas de las cuatro fases y cómo levantarlo con una orden.
  \item El cuaderno de experimentos completo y un informe técnico de diez páginas como máximo
  (plantilla en \file{docs/informe\_plantilla.md}).
  \item \textbf{Defensa oral} de diez minutos por equipo: tres de demo en vivo y el resto de
  preguntas sobre vuestro código y vuestros resultados.
\end{itemize}

\section{Por dónde empezar}

\begin{enumerate}
  \item Entrad en \url{https://github.com/kendrickcetina/dgm-arca} y pulsad
  \textbf{Use this template} para crear el repositorio de vuestro equipo. Si no habéis usado
  GitHub antes, o no sabéis cómo hacer \code{push} desde la DGX, está paso a paso en
  \file{docs/github.md}.
  \item Montad el entorno: \code{uv sync --extra train} y \code{uv run pytest}. Todo en verde
  antes de tocar nada.
  \item En la DGX: sesión de code-server desde \url{https://hpc.comillas.edu},
  \code{source smoke/dgx\_env.sh}, \code{uv run arca-check-gpu} y \code{uv run arca-smoke}.
  Guía completa, con la GPU de 16\,GB y las sesiones de 24 horas: \file{docs/dgx.md}.
  \item Escribid la propuesta (\file{docs/00\_propuesta.md}) y empezad por la sección de la
  tarea verificable. Si esa sección se escribe sola, tenéis tema.
\end{enumerate}

\textbf{Un último consejo.} Empezad por el verificador y por el conjunto dorado, no por el
modelo: lo que no se puede medir no se puede mejorar. Y haced \code{git push} al terminar cada
sesión, que la DGX no tiene copias de seguridad.

\end{document}
